\documentclass[
11pt,notheorems,hyperref={pdfauthor=whatever}
]{beamer}

% Theme
\usetheme{Madrid}
\usecolortheme{seahorse}
\beamertemplatenavigationsymbolsempty

% Packages
\usepackage[utf8]{inputenc}
\usepackage{graphicx}
\usepackage{amsmath}
\usepackage{amssymb}
\usepackage{hyperref}
\usepackage{booktabs}
\usepackage{array}
\usepackage{ragged2e}
\usepackage{tikz}
\usetikzlibrary{arrows.meta, positioning, shapes.geometric}

% Bibliography
\usepackage[backend=bibtex,style=numeric]{biblatex}
\addbibresource{../references/references.bib}

% Title page details
\title[Agentic AI]{Learning Agentic AI}

\begin{document}

% ==========================================
% SECTION 1: GUIDELINES TO PROMPT LLMS
% ==========================================

\section{Prompting LLMs}

\begin{frame}{Guidelines to Prompt LLMs for Code Generation}
	Midolo et al.\ (2026) derived and evaluated development-specific prompt optimization guidelines using an iterative, test-driven approach to automatically refine code generation prompts.

	\vspace{0.3cm}

	\textbf{Key Contributions:}
	\begin{itemize}
		\item \textbf{10 Guidelines} for prompt improvement: specifying I/O, pre/post conditions, providing examples, clarifying ambiguities.
		\item \textbf{Empirical Evaluation} with 50 practitioners rating usage and perceived usefulness.
		\item \textbf{LLM Coverage:} GPT-4o-mini, Llama 3.3 70B, Qwen2.5 72B, DeepSeek Coder V2.
	\end{itemize}

	\vspace{0.3cm}

	\textbf{Relevance to Agentic AI:} These guidelines directly inform Steps 1--3 (Goal Definition, Task Decomposition, Planning) by providing evidence-based prompt patterns that improve LLM-driven agent reliability.
\end{frame}

\begin{frame}{Example: Applying the 10 Guidelines}
	\textbf{Task:} Create a Jupyter notebook that plots a sine wave at 10\,Hz.

	\vspace{0.2cm}

	\begin{columns}[T]
		\begin{column}{0.32\textwidth}
			\textbf{\textcolor{red}{Before}}
			\begin{quote}
				\footnotesize ``Make a notebook with a sin plot.''
			\end{quote}
			\vspace{0.1cm}
			\scriptsize
			\textbf{Missing:} frequency, sample rate, axis labels, libraries.
		\end{column}
		\hfill
		\begin{column}{0.64\textwidth}
			\textbf{\textcolor{green!50!black}{After}}
			\begin{quote}
				\scriptsize
				Create a Jupyter notebook plotting a sine wave at $f=10$\,Hz.\\[1pt]
				\textbf{Requirements:} Use \texttt{numpy} and \texttt{matplotlib.pyplot}.\\[1pt]
				\textbf{Algorithm:} Sample at $f_s=1000$\,Hz over $t\in[0,1]$\,s; compute $y=\sin(2\pi f t)$.\\[1pt]
				\textbf{I/O format:} Title ``Sine Wave (10\,Hz)''; x-axis ``Time [s]'', y-axis ``Amplitude''.\\[1pt]
				\textbf{Pre/Post:} $f_s>2f$ (Nyquist); output \texttt{y} matches \texttt{t} length.\\[1pt]
				\textbf{Errors:} Raise \texttt{ValueError} if $f_s\leq0$ or $f\leq0$.\\[1pt]
				\textbf{Example:} For $f=1$, $t=[0,0.1,\dots,0.9]$, expect $y[0]=0$, $y[2]\approx0.588$.
			\end{quote}
		\end{column}
	\end{columns}

	\vspace{0.25cm}

	\tiny
	\textbf{Guidelines applied:} Requirements \checkmark\; Pre/Post-conditions \checkmark\; I/O Format \checkmark\; Exceptions \checkmark\; Algorithmic Details \checkmark\; Variable Naming \checkmark\; Clarity \checkmark\; Examples \checkmark\; Assertive Language \checkmark
\end{frame}

\begin{frame}{Midolo et al.\ (2026): 10 Prompt Guidelines}
	\centering
	\scriptsize
	\begin{tabular}{p{3.2cm}p{6.5cm}r}
		\toprule
		\textbf{Guideline} & \textbf{Description} & \textbf{Applied} \\
		\midrule
		Requirements & Make required packages/libraries explicit; explain their purpose. & 19\% \\
		Pre-conditions & State conditions that must hold before execution. & 7\% \\
		Post-conditions & Specify expected guarantees after execution. & 23\% \\
		I/O format & Clarify input/output types, shapes, and special cases. & 44\% \\
		Exceptions & Explicitly define exceptions to raise and error messages. & 12\% \\
		Algorithmic Details & Detail formulas, definitions, and complex logic. & 57\% \\
		Variable Mentioning & Use consistent terminology for variables. & 3\% \\
		Unclear Conditions & Avoid ``otherwise''; describe both branches explicitly. & 1\% \\
		More Examples & Add doctests/run examples. & 24\% \\
		Assertive Language & Use ``must'' instead of ``should''; avoid ambiguity. & 9\% \\
		\bottomrule
	\end{tabular}

	\vspace{0.25cm}

	\footnotesize
	\textbf{Takeaway:} Algorithmic details and I/O format are the most frequently needed refinements; linguistic fixes (variable naming, unclear conditions) are rarer but still valuable.
\end{frame}

% ==========================================
% SECTION 2: FUNCTION-LEVEL CODE GENERATION
% ==========================================

\section{Function-Level Code Generation}

\begin{frame}{CodePromptEval: Dataset \& Methodology}
	Khojah et al.\ (2025) introduced \textbf{CodePromptEval}, a full factorial dataset to evaluate how prompt techniques---and their interactions---impact function-level code generation.

	\vspace{0.25cm}

	\begin{columns}[T]
		\begin{column}{0.55\textwidth}
			\textbf{Experimental Design}
			\begin{itemize}
				\item \textbf{221 tasks} from CoderEval $\times$ \textbf{32 variations} = \textbf{7,072 prompts}
				\item \textbf{5 techniques:} few-shot, CoT, persona, signature, packages
				\item \textbf{3 LLMs:} GPT-4o, Llama3-70B, Mistral-22B (3 runs each, temp=0.2)
			\end{itemize}
		\end{column}
		\hfill
		\begin{column}{0.43\textwidth}
			\textbf{Evaluation}
			\begin{itemize}
				\item \textbf{Correctness:} Pass@1
				\item \textbf{Similarity:} CrystalBLEU
				\item \textbf{Quality:} code smells + complexity
			\end{itemize}
		\end{column}
	\end{columns}

	\vspace{0.25cm}

	\textbf{Relevance:} Systematic benchmark for evaluating prompt engineering strategies in code generation agents.
\end{frame}

\begin{frame}{Key Findings: The Correctness--Quality Trade-off}
	\begin{columns}[T]
		\begin{column}{0.48\textwidth}
			\textbf{Correctness (RQ2.1)}
			\begin{itemize}
				\item \textbf{Signature + few-shot} = best Pass@1 (57.5\% for GPT-4o)
				\item Best vs.\ worst combo: only \textbf{10 pp gap}
				\item Combining techniques often \textbf{cancels out} benefits
			\end{itemize}
		\end{column}
		\hfill
		\begin{column}{0.48\textwidth}
			\textbf{Quality (RQ2.3)}
			\begin{itemize}
				\item CoT, persona, packages $\rightarrow$ \textbf{fewer code smells}
				\item Few-shot + signature $\rightarrow$ \textbf{more code smells} (71\% lack docstrings)
				\item Llama3 produces simpler code than GPT-4o
			\end{itemize}
		\end{column}
	\end{columns}

	\vspace{0.3cm}

	\begin{block}{\centering Critical Trade-off (L3)}
		\centering
		Correctness techniques (signature, few-shot) worsen quality.\\
		Quality techniques (persona, CoT) reduce pass rates.\\
		\textbf{Agents must optimize for the right metric per task.}
	\end{block}

	\vspace{0.1cm}

	\textbf{Relevance to Agentic AI:} Blindly stacking prompt techniques is counterproductive.
\end{frame}

% ==========================================
% SECTION 3: SOFTWARE PRACTITIONERS
% ==========================================

\section{Software Practitioners}

\begin{frame}{How Practitioners Actually Use GenAI for SE}
	Otten et al.\ (2025) surveyed \textbf{72 software practitioners} across 6 SE tasks to characterize real-world GenAI usage, prompting strategies, and reliability perceptions.

	\vspace{0.25cm}

	\begin{columns}[T]
		\begin{column}{0.48\textwidth}
			\textbf{Tool Adoption}
			\begin{itemize}
				\item \textbf{ChatGPT} dominates (83.1\%), Copilot second (43.7\%)
				\item Web + IDE equally preferred (63.9\% vs.\ 62.5\%)
				\item Daily use correlates with task coverage ($\rho = 0.325$)
			\end{itemize}
		\end{column}
		\hfill
		\begin{column}{0.48\textwidth}
			\textbf{Workflow Impact}
			\begin{itemize}
				\item 77.5\% report GenAI changed their approach
				\item Replaced web searches and boilerplate generation
				\item Shifted focus to higher-level design
			\end{itemize}
		\end{column}
	\end{columns}

	\vspace{0.25cm}

	\textbf{Relevance:} Understanding practitioner workflows informs how agents should integrate into existing development practices.
\end{frame}

\begin{frame}{Prompting in Practice: Strategies \& Reliability}
	\begin{columns}[T]
		\begin{column}{0.48\textwidth}
			\textbf{Conversation Strategies}
			\begin{itemize}
				\item \textbf{Iterative multi-turn} preferred over single-shot
				\item Incremental refinement + feedback loop = most common
				\item Proficient users require \textbf{10+ exchanges} for complex tasks
			\end{itemize}

			\vspace{0.2cm}

			\textbf{Awareness vs.\ Adoption}
			\begin{itemize}
				\item 74--94\% familiar, but $<$50\% regularly use them
				\item Intermediates (6--10 yr) = most active adopters
			\end{itemize}
		\end{column}
		\hfill
		\begin{column}{0.48\textwidth}
			\textbf{Reliability Perception}
			\begin{itemize}
				\item \textbf{Most reliable:} documentation (avg.\ 3.06/4)
				\item \textbf{Least reliable:} debugging, complex algorithms (avg.\ 2.34--2.38/4)
				\item \#1 issue: ``bug fixes miss root causes''
			\end{itemize}
		\end{column}
	\end{columns}

	\vspace{0.3cm}

	\textbf{Relevance to Agentic AI:} Agents should emulate iterative refinement, embed techniques into interfaces, and prioritize reliability for debugging.
\end{frame}

% ==========================================
% REFERENCES
% ==========================================

\begin{frame}[allowframebreaks]{References}
	\nocite{*}
	\printbibliography
\end{frame}

\end{document}
